\section{The classical single-exponential translation}\label{app:classical-translation}
We include the size analysis for the classical expressive equivalence
\citep{tarski-givant1987,maddux2006}. The relevant point is that successive
quantifier eliminations need not create a tower of exponentials: an
already translated subformula is treated as one leaf of the next Boolean
expansion, without expanding its algebraic representation again.

\begin{proposition}[Classical translation]\label{prop:upper-translation}
Every \FOthree\ formula $\varphi$ over binary relations with at most two
free variables has an equivalent \CoR\ term of size $2^{O(|\varphi|)}$
on all nonempty structures, constructible in time $2^{O(|\varphi|)}$.
Here size counts syntax-tree nodes.
\end{proposition}
\begin{proof}
We use the binary-assignment convention of \S\ref{sec:prelim-logic}:
represent each subformula with at most two free variables on an
ordered pair containing them. An unused coordinate is a dummy, so
moving the representation to another suitable pair requires only
converse (if the coordinates are exchanged). On each structure, the
representation of a sentence has value $\zero$ or $\one$, according
to its truth value on that structure.

The atomic translations are those of \S\ref{sec:FOtwo-translation},
with the two relevant variables renamed or reordered: distinct-variable
atoms use the input relation or its converse, reflexive atoms use a
diagonal test followed by a projection, and equality uses $\id$ or
$\one$. All have constant size. Boolean operations on a fixed
variable pair become intersection, union, and complement.

Process quantified subformulas from the inside out. Consider
$\exists w\,\psi$, and let $u,v$ be the other two variable names, in
that order. Its free variables are contained in $\{u,v\}$.
In the body $\psi$, regard each atom and each maximal quantified
subformula as one leaf of a Boolean formula, called its \defn{shell}.
Each leaf has at most two free variables: this is immediate for binary
atoms, and a quantifier binds one of the three available names.
The quantified leaves have already been translated. Crucially, the
shell expansion treats each of their translations as an opaque term.

If the shell has $r$ leaf occurrences, a truth-table DNF has at most
$2^r$ clauses, each using at most $r$ literals. Repeated or logically
dependent leaves cause no problem; the propositional identity remains
valid after substituting their actual interpretations. In each clause,
assign its literals to three groups according to the free-variable pair:
$(u,v)$, $(u,w)$, or $(w,v)$. A literal with only one free variable is
assigned to any group containing it; a closed literal can be put in the
first group. Use converse to orient each translated literal for its
chosen pair, and use complement for a negative literal.
Writing the conjunctions of the groups as $A(u,v)$, $B(u,w)$, and
$C(w,v)$, apply the identity~\eqref{eq:prelim-elimination}:
\[
 \exists w\,(A(u,v)\land B(u,w)\land C(w,v))
       \quad\longmapsto\quad A\cap(B;C).
\]
An empty group is represented by $\one$. This identity also covers
vacuous quantification because domains are nonempty.
Take the union of the translated clauses. Universal quantification uses
$\forall w\,\psi\equiv\neg\exists w\,\neg\psi$, which introduces no
subformula duplication. At the root, any remaining Boolean shell with
at most two free variables is translated directly, after orienting its
leaves to a common ordered pair. This proves correctness on all nonempty
structures by induction.

For the size bound, cut the original syntax tree at
quantifiers into the shells just used. Count each child quantified
subformula as one leaf of its parent's shell. The sum of the shell sizes,
including these boundary leaves, is $O(|\varphi|)$: each original node is counted
once, with only one extra boundary occurrence per quantified subformula.
A shell of size $s$ creates at most $2^{O(s)}$ local nodes and at most
$2^{O(s)}$ copies of each child translation. Along any path from an
original leaf or a newly introduced local node to the final output,
the resulting duplication factor is therefore at most
\[
 \prod_{\text{shells on the path}}2^{O(s)}
       \le 2^{O(|\varphi|)}.
\]
There are $O(|\varphi|)$ original nodes and shell locations, and the local
exponential creation factor is included in this same bound. Summing
their contributions gives $|\varphi|\,2^{O(|\varphi|)}=2^{O(|\varphi|)}$ nodes.
No step expands the internal Boolean structure of a translated child,
which is what prevents an iterated-exponential recurrence.
Enumerating the shell assignments and writing the resulting term
takes a polynomial factor in $|\varphi|$ times the same exponential
bound, hence also $2^{O(|\varphi|)}$ time.
\end{proof}
